\documentclass[11pt]{article}
\usepackage{amsmath,amssymb,amsthm}

\newtheorem{theorem}{Theorem}
\newtheorem{lemma}[theorem]{Lemma}
\newtheorem{corollary}[theorem]{Corollary}

\newcommand{\Up}{U_+}

\title{Tail-record renewal at the deficit-four horizon}
\author{}
\date{27 August 2026}

\begin{document}
\maketitle

Let
\[
  x_{t+1}=\Up(x_t):=\frac{3x_t+1}{2^{a_t}},
  \qquad
  a_t:=v_2(3x_t+1),
  \qquad
  A_t:=\sum_{h=0}^{t-1}a_h,
  \tag{1}
\]
be an infinite accelerated orbit of positive odd integers.  Define its
coefficient sequence and decoder depth by
\[
  Q_t:=\frac{3^t}{2^{A_t}},
  \qquad
  b_t:=v_3(x_t+1).
  \tag{2}
\]
Assume throughout that $Q_t\to\infty$.  A time $j$ is a tail record if
\[
  Q_j<Q_t\qquad(t>j).
  \tag{3}
\]
The tail-record times form an infinite sequence
\[
  j_0<j_1<j_2<\cdots.
  \tag{4}
\]

For a positive word $w=(c_0,\ldots,c_{r-1})$, put
\[
\begin{aligned}
  A_0(w)&:=0,
  &A_h(w)&:=\sum_{t=0}^{h-1}c_t,\\
  C_0(w)&:=0,
  &C_{h+1}(w)&:=3C_h(w)+2^{A_h(w)}.
\end{aligned}
\tag{5}
\]
and
\[
  D(w):=C_r(w)+2^{A_r(w)}.
  \tag{6}
\]
The affine identity is
\[
  2^{A_r(w)}(x_{j+r}+1)=3^r x_j+D(w)
  \tag{7}
\]
whenever $w=(a_j,\ldots,a_{j+r-1})$.

\begin{lemma}[geometry of consecutive tail records]
Put $r_i:=j_{i+1}-j_i$ and
\[
  B_h:=A_{j_i+h}-A_{j_i},
  \qquad
  q_h:=\frac{Q_{j_i+h}}{Q_{j_i}}
       =\frac{3^h}{2^{B_h}}.
  \tag{8}
\]
Then
\[
  1<q_{r_i}<q_h
  \qquad(1\leq h<r_i),
  \qquad
  a_{j_i}=a_{j_{i+1}}=1.
  \tag{9}
\]
If $r_i>1$, then $a_{j_{i+1}-1}\geq2$.

The consecutive-record gaps of length less than six are exactly
\[
  \begin{array}{c|c}
    r_i& (a_{j_i},\ldots,a_{j_{i+1}-1})\\ \hline
    1&1\\
    2&12\\
    4&1113\text{ or }1122.
  \end{array}
  \tag{10}
\]
There are no consecutive-record gaps of length $3$ or $5$.
\end{lemma}

\begin{proof}
The values $Q_t$ are pairwise distinct, since an equality $Q_u=Q_v$ with
$u<v$ would give a positive power of $3$ equal to a power of $2$.
Because $Q_t\to\infty$, every suffix has a unique minimum.  The next tail
record after $j_i$ is the unique minimum of $\{Q_t:t>j_i\}$.  This proves
the first two inequalities in (9).  At every tail record $j$, one has
\[
  \frac{Q_{j+1}}{Q_j}=\frac3{2^{a_j}}>1,
\]
so $a_j=1$.  If $r_i>1$, then $j_{i+1}-1$ is not a tail record, and (9)
gives
\[
  \frac{Q_{j_{i+1}}}{Q_{j_{i+1}-1}}
  =\frac3{2^{a_{j_{i+1}-1}}}<1.
\]
Hence the last exponent is at least two.

It remains to prove (10).  The endpoint condition $q_r<q_1=3/2$ and the
strict barrier $q_r>1$ give
\[
  1<\frac{3^r}{2^{B_r}}<\frac32.
  \tag{11}
\]
For $r=2$, this forces $B_2=3$, hence the word is $12$.  For $r=3$,
(11) forces $B_3\geq5$, while $q_3>1$ forces $B_3\leq4$; hence no word
exists.  For $r=4$, equation (11) forces $B_4=6$.  The prefix barriers
leave only
\[
  1113,\qquad1122,\qquad1212.
\]
The last word is impossible because its second coefficient is $9/8$, less
than its endpoint coefficient $81/64$.  The first two words satisfy (9).
For $r=5$, equation (11) forces $B_5\geq8$, whereas $q_5>1$ forces
$B_5\leq7$.  The case $r=1$ was already proved.  This establishes the
table.
\end{proof}

Let
\[
  \mathcal S_4=
  \{\varnothing,12,112,1112,1113,1122,1212,11113\}.
  \tag{12}
\]
A stem $s\in\mathcal S_4$ matches a positive odd integer $x$ if it is the
initial accelerated exponent word of $x$.  Define
\[
  R_4(x):=
  \max_{\substack{s\in\mathcal S_4\\s\text{ matches }x}}
  \left(
    |s|+v_2\bigl(\Up^{|s|}(x)+1\bigr)-1
  \right).
  \tag{13}
\]
The empty stem makes the maximum well defined.

\begin{lemma}[exact exponent-one tower]
For every positive odd $y$, the exact initial run of exponent-one steps
has length
\[
  v_2(y+1)-1.
  \tag{14}
\]
Consequently, if a matching stem $s$ contributes the value $R$ to the
maximum in (13), then the exponent prefix of length $R$ is
\[
  s1^{R-|s|}.
  \tag{15}
\]
\end{lemma}

\begin{proof}
Put $\nu=v_2(y+1)$.  If $\nu\geq2$, then $y\equiv-1\pmod4$, so
$v_2(3y+1)=1$, and
\[
  \Up(y)+1=\frac{3(y+1)}2.
  \tag{16}
\]
Thus one exponent-one step lowers $\nu$ by exactly one.  When $\nu=1$,
one has $y\equiv1\pmod4$, so the next exponent is at least two.  This
proves both assertions.
\end{proof}

\begin{theorem}[deficit-four record-renewal dichotomy]
For each $i$, put
\[
  r_i:=j_{i+1}-j_i,
  \qquad
  R_i:=R_4(x_{j_i}).
  \tag{17}
\]
Exactly one of the following alternatives holds.

\begin{enumerate}
  \item If $r_i>R_i$, then
  \[
    r_i\geq6,
    \qquad
    b_{j_{i+1}}\leq r_i-5.
    \tag{18}
  \]

  \item If $r_i\leq R_i$, then the record-gap word $p_i$ belongs to
  \[
    \mathcal A_4:=\{1,12,1113,1122\},
    \tag{19}
  \]
  and
  \[
    R_{i+1}\geq R_i-r_i.
    \tag{20}
  \]
  Its decoder transport is exact:
  \[
    \begin{array}{c|c|c|c}
      p_i&r_i&A_{r_i}(p_i)&b_{j_{i+1}}\\ \hline
      1&1&1&b_{j_i}+1\\
      12&2&3&0\\
      1113&4&6&1\\
      1122&4&6&0.
    \end{array}
    \tag{21}
  \]
\end{enumerate}
\end{theorem}

\begin{proof}
Fix $i$ and suppose first that $r_i\leq R_i$.  Choose a stem $s$ that
attains the maximum $R_i$ in (13).  By Lemma~2, the prefix through time
$R_i$ is $s$ followed by exponent-one steps.  Since $j_{i+1}$ is the
unique minimum of the coefficient sequence after $j_i$, its relative
coefficient is smaller than that at every other positive time through
$R_i$.

The required finite comparison is
\[
  \begin{array}{c|c|c}
    s& p_i&\text{stem remaining after }p_i\\ \hline
    \varnothing&1&\varnothing\\
    12&12&\varnothing\\
    112&1&12\\
    1112&1&112\\
    1113&1113&\varnothing\\
    1122&1122&\varnothing\\
    1212&12&12\\
    11113&1&1113.
  \end{array}
  \tag{22}
\]
Indeed, the relative coefficient vectors through the nonempty stems are
\[
\begin{aligned}
12&:\quad (3/2,9/8),\\
112&:\quad (3/2,9/4,27/16),\\
1112&:\quad (3/2,9/4,27/8,81/32),\\
1113&:\quad (3/2,9/4,27/8,81/64),\\
1122&:\quad (3/2,9/4,27/16,81/64),\\
1212&:\quad (3/2,9/8,27/16,81/64),\\
11113&:\quad (3/2,9/4,27/8,81/16,243/128).
\end{aligned}
\tag{23}
\]
The unique minimum in each row occurs at the boundary displayed in (22).
Every appended exponent-one step multiplies the coefficient by $3/2$, so
it creates no new minimum.  This proves (19).  The remaining word through
time $R_i$ is again a stem from $\mathcal S_4$ followed by its exact
exponent-one tower.  Applied at $x_{j_{i+1}}$, definition (13) therefore
gives (20).

Direct use of (5)--(6) gives
\[
  \begin{array}{c|rrrr}
    p&1&12&1113&1122\\ \hline
    D(p)&3&13&129&137\\
    v_3(D(p))&1&0&1&0.
  \end{array}
  \tag{24}
\]
For $p=1$, equation (7) reduces to
\[
  2(x_{j_i+1}+1)=3(x_{j_i}+1),
\]
which gives the first decoder rule in (21).  For the other three words,
$v_3(D(p))<|p|$, while the first term on the right of (7) has valuation at
least $|p|$.  Hence the valuation of the sum is $v_3(D(p))$, proving the
remaining rules.

Now suppose that $r_i>R_i$.  Lemma~1 and the definition (13) show that none
of the gaps in (10) can occur: each of their words is a matching stem and
therefore gives $R_i\geq r_i$.  Thus $r_i\geq6$.  The block from
$x_{j_i}$ is a strict barrier because $j_i$ is a tail record.  The proved
fixed-start deficit-four horizon theorem, applied with $r_i\geq4$ and
$r_i>R_i$, gives the second assertion in (18).
\end{proof}

\begin{corollary}[record-renewal trichotomy]
Every orbit satisfying the hypotheses has at least one of the following
three properties.

\begin{enumerate}
  \item Infinitely many consecutive tail-record transitions satisfy
  $r_i>R_i$.

  \item From some tail record onward every transition lies inside its
  current horizon, and the resulting word over $\mathcal A_4$ contains
  arbitrarily long consecutive runs of the one-letter atom $1$.

  \item From some tail record onward every transition lies inside its
  current horizon, and the decoder depth $b_t$ is bounded on that tail.
\end{enumerate}
In the third alternative, if every run of the atom $1$ has length at most
$M$, then one may take
\[
  b_t\leq M+4
  \tag{25}
\]
throughout a sufficiently late tail.
\end{corollary}

\begin{proof}
If the first alternative fails, Theorem~3 shows that all sufficiently late
record gaps belong to $\mathcal A_4$.  The three nontrivial atoms occur
infinitely often.  Otherwise the accelerated exponent word would
eventually be all ones, contradicting the finite exact run length in
Lemma~2.

If the intervening runs of the atom $1$ are unbounded, the second
alternative holds.  Otherwise let $M$ bound their lengths.  Equation (21)
shows that each nontrivial atom resets the decoder depth to zero or one,
whereas each atom $1$ raises it by one.  Inside a nontrivial atom there are
at most three initial exponent-one steps before the displayed reset.
Consequently $b_t\leq M+4$ on the resulting tail, proving the third
alternative and (25).
\end{proof}

\begin{theorem}[record carry and explicit coalescing ancestor]
Fix consecutive tail records $j_i<j_{i+1}$, put
\[
  r:=j_{i+1}-j_i>1,
  \qquad
  x:=x_{j_i},
  \qquad
  y:=x_{j_{i+1}},
  \qquad
  B:=A_{j_{i+1}}-A_{j_i},
  \tag{26}
\]
and let $w=(a_{j_i},\ldots,a_{j_{i+1}-1})$.  With
$\alpha:=\log_2 3$ and $\{u\}:=u-\lfloor u\rfloor$, one has
\[
  B=\lfloor r\alpha\rfloor,
  \qquad
  0<\{r\alpha\}<\alpha-1,
  \qquad
  \frac{3^r}{2^B}=2^{\{r\alpha\}}\in(1,3/2).
  \tag{27}
\]
Moreover,
\[
  \frac{C(w)}{2^B}
  =\frac13\sum_{h=0}^{r-1}\frac{q_r}{q_h}
  <\frac{r+1/2}{3},
  \tag{28}
\]
where the $q_h$ are defined in (8), and consequently
\[
  \frac{3^r}{2^B}x
  <y
  <\frac{3^r}{2^B}x+\frac{r+1/2}{3}
  <\frac32x+\frac{r+1/2}{3}.
  \tag{29}
\]

Put $b:=v_3(y+1)$.  Then
\[
  b\geq1
  \quad\Longleftrightarrow\quad
  a_{j_{i+1}-1}\text{ is odd}.
  \tag{30}
\]
If these equivalent conditions hold, the positive odd integer
\[
  z:=2^b\frac{y+1}{3^b}-1
  \tag{31}
\]
satisfies
\[
  \Up^b(z)=y,
  \tag{32}
\]
and all $b$ exponents on this coalescing path are one.  It is smaller than
the record start $x$ exactly when
\[
  \left(\frac23\right)^b(y+1)<x+1.
  \tag{33}
\]
In particular, either of the following explicit inequalities is sufficient:
\[
  \left(\frac23\right)^b
  \left(
    \frac32x+\frac{r+1/2}{3}+1
  \right)<x+1,
  \tag{34}
\]
or
\[
  b\geq2
  \quad\text{and}\quad
  4r<9x+13.
  \tag{35}
\]
\end{theorem}

\begin{proof}
Lemma~1 gives
\[
  1<\frac{3^r}{2^B}<\frac32.
\]
Taking base-two logarithms gives
\[
  0<r\alpha-B<\alpha-1<1.
\]
Since $B$ is an integer, this proves (27).

Expansion of the carry recurrence in (5) gives
\[
  C(w)=\sum_{h=0}^{r-1}3^{r-1-h}2^{B_h}.
\]
After division by $2^B$, each summand becomes
\[
  \frac13\frac{3^{r-h}}{2^{B-B_h}}
  =\frac13\frac{q_r}{q_h}.
\]
For $h=0$, Lemma~1 gives $q_r/q_0=q_r<3/2$.  For
$1\leq h<r$, it gives $q_r/q_h<1$.  Summation proves (28).
The affine orbit identity
\[
  y=\frac{3^r}{2^B}x+\frac{C(w)}{2^B}
\]
then proves (29).

Write $a:=a_{j_{i+1}-1}$.  The final orbit step gives
\[
  2^a y\equiv1\pmod3.
\]
Thus $y\equiv(-1)^a\pmod3$.  Hence $3\mid y+1$ exactly when $a$ is odd,
which proves (30).

Suppose $b\geq1$.  Formula (31) defines a positive odd integer.  For
$0\leq h\leq b$, put
\[
  z_h+1:=2^{b-h}\frac{y+1}{3^{b-h}}.
  \tag{36}
\]
Every $z_h$ is a positive odd integer, $z_0=z$, and $z_b=y$.  Moreover,
\[
  z_{h+1}+1=\frac{3(z_h+1)}2.
  \tag{37}
\]
Since $j_{i+1}$ is a tail record, Lemma~1 gives
$a_{j_{i+1}}=1$, and therefore $v_2(y+1)\geq2$.  It follows from (36)
that $v_2(z_h+1)\geq2$ for $h<b$.  Lemma~2 and (37) now show that each
step from $z_h$ to $z_{h+1}$ has exact exponent one.  This proves (32).

Equation (31) gives
\[
  z+1=\left(\frac23\right)^b(y+1),
\]
which proves (33).  Combining this identity with (29) proves (34).  If
$b\geq2$, the left side of (34) is at most
\[
  \frac49
  \left(
    \frac32x+\frac{r+1/2}{3}+1
  \right)
  =
  \frac23x+\frac{4r+14}{27}.
\]
This is less than $x+1$ whenever $4r<9x+13$, proving (35).
\end{proof}

\begin{theorem}[weighted monotonicity of decoder ancestors]
Suppose that the transition from $j_i$ to $j_{i+1}$ lies inside the
current deficit-four horizon.  Put
\[
  \mathcal W_i:=
  2^{b_{j_i}}\frac{x_{j_i}+1}{3^{b_{j_i}}},
  \qquad
  \kappa(b):=
  \begin{cases}
    1,&b=0,\\
    32/27,&b\geq1.
  \end{cases}
  \tag{38}
\]
Then $\mathcal W_i-1$ is a positive all-one coalescing ancestor of
$x_{j_i}$, and the four interior record atoms act by
\[
\begin{array}{c|c|c}
  p_i&b_{j_{i+1}}&\mathcal W_{i+1}\\ \hline
  1&b_{j_i}+1&\mathcal W_i\\[2pt]
  12&0&
  \displaystyle
  \frac98\left(\frac32\right)^{b_{j_i}}\mathcal W_i+\frac12\\[7pt]
  1113&1&
  \displaystyle
  \frac{27}{32}\left(\frac32\right)^{b_{j_i}}\mathcal W_i+\frac12\\[7pt]
  1122&0&
  \displaystyle
  \frac{81}{64}\left(\frac32\right)^{b_{j_i}}\mathcal W_i+\frac78.
\end{array}
\tag{39}
\]
In every case,
\[
  \boxed{
  \kappa(b_{j_{i+1}})\mathcal W_{i+1}
  \geq
  \kappa(b_{j_i})\mathcal W_i.}
  \tag{40}
\]
Consequently, along any finite or infinite chain of interior record
renewals,
\[
  \mathcal W_k\geq\frac{27}{32}\mathcal W_i
  \qquad(k\geq i).
  \tag{41}
\]
In particular, every run of the one-letter atom $1$ leaves the canonical
coalescing ancestor exactly unchanged; unbounded exponent-one towers do not
produce cumulative ancestor descent.
\end{theorem}

\begin{proof}
For any positive odd state $v$, let
\[
  b:=v_3(v+1),
  \qquad
  \mathcal W(v):=2^b\frac{v+1}{3^b}.
\]
The construction in the proof of Theorem~5 shows that
$\mathcal W(v)-1$ reaches $v$ through exactly $b$ exponent-one steps.
This proves the first assertion.

Write $Y:=x_{j_i}+1$, $Y':=x_{j_{i+1}}+1$, and
$b:=b_{j_i}$.  Equations (7) and (24) give, for the four atoms,
\[
\begin{array}{c|c}
  p_i&\text{block equation}\\ \hline
  1&2Y'=3Y\\
  12&8Y'=9Y+4\\
  1113&64Y'=81Y+48\\
  1122&64Y'=81Y+56.
\end{array}
\tag{42}
\]
Since
\[
  Y=\left(\frac32\right)^b\mathcal W_i,
  \tag{43}
\]
substitution into (42), together with the decoder endpoints in (21),
proves every formula in (39).  In particular, the factor $3/2$ in the
one-letter block raises the decoder depth at the same time as it changes
$Y$, so $\mathcal W_{i+1}=\mathcal W_i$.

It remains to verify (40).  For the atom $1$, the weight $\kappa$ increases
from $1$ to $32/27$ when $b=0$ and is otherwise unchanged.  For the atom
$12$, the leading multiplier in (39) is at least $9/8$ when $b=0$ and at
least $27/16>32/27$ when $b\geq1$.  For the atom $1113$, multiplication
by the terminal weight gives the leading factor
\[
  \frac{32}{27}\frac{27}{32}
  \left(\frac32\right)^b
  =
  \left(\frac32\right)^b.
\]
This equals the initial weight when $b=0$ and exceeds $32/27$ when
$b\geq1$.  Finally, the leading multiplier for $1122$ is at least
$81/64>1$ when $b=0$ and at least
$243/128>32/27$ when $b\geq1$.  All additive terms in (39) are positive,
so (40) follows.  Iteration and
$1\leq\kappa(b)\leq32/27$ give (41).
\end{proof}

\begin{theorem}[escape of canonical decoder ancestors]
Let $(x_t)_{t\geq0}$ be an injective infinite accelerated orbit.  For every
$t$, put
\[
  b_t:=v_3(x_t+1),
  \qquad
  \mathcal W_t:=2^{b_t}\frac{x_t+1}{3^{b_t}}.
  \tag{44}
\]
Then
\[
  \boxed{\mathcal W_t\longrightarrow\infty.}
  \tag{45}
\]
More precisely, for every positive integer $M$,
\[
  \#\{t\geq0:\mathcal W_t=M\}
  \leq v_2(M).
  \tag{46}
\]
The left side is zero when $M$ is odd.

Suppose additionally that $x_0=n$ is the least positive odd seed with an
unbounded accelerated orbit.  Then
\[
  \mathcal W_t\geq n+1
  \qquad(t\geq0),
  \tag{47}
\]
with equality only during the exact initial exponent-one tower of $n$.
In particular, after that finite tower every canonical decoder ancestor is
strictly larger than $n$.
\end{theorem}

\begin{proof}
Every $\mathcal W_t$ is even.  Indeed, if $b_t=0$, then
$\mathcal W_t=x_t+1$ is even, while the assertion is immediate when
$b_t\geq1$.

Fix an even integer $M$ and put $z:=M-1$.  If $\mathcal W_t=M$, the
construction in Theorem~5 gives
\[
  \Up^{b_t}(z)=x_t,
  \tag{48}
\]
and all $b_t$ exponents on this path are one.  Lemma~2 therefore gives
\[
  0\leq b_t\leq v_2(M)-1.
  \tag{49}
\]
For a fixed value of $b_t$, equation (48) determines $x_t$.  Injectivity of
the orbit makes the values $x_t$ pairwise distinct, so each integer in the
range (49) occurs for at most one $t$.  This proves (46).

For any fixed $H$, only finitely many even integers $M\leq H$ exist, and
each occurs only finitely often by (46).  Hence $\mathcal W_t>H$ for all
sufficiently large $t$.  Since $H$ is arbitrary, (45) follows.

Finally, the orbit of $\mathcal W_t-1$ reaches the unbounded state $x_t$.
It is therefore unbounded, so global leastness of $n$ gives (47).  If
we take $t=0$ in (47), the defining formula (44) first shows that
$b_0=0$ and $\mathcal W_0=n+1$.  If equality holds at any $t$, then (48)
becomes
\[
  \Up^{b_t}(n)=x_t=\Up^t(n).
\]
Injectivity gives $b_t=t$.  The coalescing path in (48) then shows that the
first $t$ exponents of $n$ are all one.  Conversely, every state in the
exact initial exponent-one tower has canonical ancestor $n$.  Lemma~2
shows that this tower is finite, completing the proof.
\end{proof}

\begin{theorem}[two-prime stationary budget at coefficient records]
Assume in addition that the normalized height has a finite limit
\[
  \mathcal K:=\lim_{t\to\infty}\frac{x_t}{Q_t}<\infty.
  \tag{50}
\]
Put
\[
  \nu_t:=v_2(x_t+1),
  \qquad
  G_t:=\frac{x_t+1}{Q_t},
  \qquad
  \gamma:=\log_3\frac32.
  \tag{51}
\]
Then, for every $t$,
\[
  \boxed{
  2^{\nu_t}3^{b_t}
  \leq x_t+1
  =Q_tG_t
  \leq \mathcal KQ_t.}
  \tag{52}
\]
Consequently,
\[
  \boxed{
  \mathcal W_t
  \geq
  G_0\left(\frac2{\mathcal K}\right)^\gamma
  Q_t^{\log_3 2}.}
  \tag{53}
\]
In particular, $Q_t\to\infty$ quantitatively forces
$\mathcal W_t\to\infty$ and
\[
  \frac{Q_t}{(3/2)^{b_t}}\longrightarrow\infty.
  \tag{54}
\]

Let
\[
  \alpha:=\log_2 3,
  \qquad
  \ell_i:=\nu_{j_i}-1.
  \tag{55}
\]
Thus $\ell_i$ is the exact initial exponent-one run at the $i$th tail
record.  Then
\[
  \boxed{
  \ell_i+1+\alpha b_{j_i}
  \leq
  (\alpha-1)i+\log_2(\mathcal KQ_{j_0}).}
  \tag{56}
\]
Since $\ell_i\geq1$, this gives the explicit record-index ceiling
\[
  \boxed{
  b_{j_i}
  \leq
  \frac{\alpha-1}{\alpha},i
  +\frac{\log_2(\mathcal KQ_{j_0})-2}{\alpha}.}
  \tag{57}
\]
Hence $i-b_{j_i}\to\infty$.

Finally, let $L_i$ be the number of consecutive one-letter record atoms
beginning at $j_i$.  This number is finite and
\[
  \boxed{
  L_i+2+\alpha b_{j_i}
  \leq
  (\alpha-1)i+\log_2(\mathcal KQ_{j_0}).}
  \tag{58}
\]
Thus the two interior mechanisms consume the same stationary budget:
decoder depth has asymptotic record-index slope at most
$(\alpha-1)/\alpha$, while an exponent-one record tower has slope at most
$\alpha-1$.
\end{theorem}

\begin{proof}
The coprime powers $2^{\nu_t}$ and $3^{b_t}$ both divide $x_t+1$, so
\[
  2^{\nu_t}3^{b_t}\mid x_t+1.
  \tag{59}
\]
To prove the upper bound in (52), first note that
\[
  G_{t+1}-G_t
  =\frac{2^{a_t}-2}{3Q_t}\geq0.
  \tag{60}
\]
Moreover, $G_t=x_t/Q_t+1/Q_t\to\mathcal K$.  Hence
$G_t\leq\mathcal K$, and (52) follows from (59).

Since $\nu_t\geq1$, equation (52) gives
\[
  3^{b_t}\leq\frac{\mathcal KQ_t}{2}.
  \tag{61}
\]
Raising (61) to the power
$\gamma=\log_3(3/2)$ yields
\[
  \left(\frac32\right)^{b_t}
  \leq
  \left(\frac{\mathcal KQ_t}{2}\right)^\gamma.
  \tag{62}
\]
On the other hand, (44) and (51) give
\[
  \mathcal W_t
  =\frac{Q_tG_t}{(3/2)^{b_t}}.
  \tag{63}
\]
The monotonicity in (60) gives $G_t\geq G_0$.  Substitution of (62) into
(63), together with $1-\gamma=\log_3 2$, proves (53).  Equations (53) and
(63) prove (54).

By Lemma~1,
\[
  \frac{Q_{j_{i+1}}}{Q_{j_i}}\leq\frac32,
\]
and therefore
\[
  Q_{j_i}\leq Q_{j_0}\left(\frac32\right)^i.
  \tag{64}
\]
Apply (52) at $j_i$ and take binary logarithms.  Since
$\nu_{j_i}=\ell_i+1$, equation (64) gives (56).  Every tail record begins
with an exponent-one step, so $\ell_i\geq1$; this proves (57), and its
linear coefficient is strictly less than one.

It remains to prove (58).  An infinite sequence of consecutive one-letter
atoms would give an infinite exponent-one run, contrary to Lemma~2.  Thus
$L_i$ is finite.  The $L_i$ one-letter gaps are followed by another record
gap, whose first exponent is again one by Lemma~1.  Hence
$\ell_i\geq L_i+1$.  Substitution in (56) proves (58).
\end{proof}

\end{document}
